% Part: sets-functions-relations
% Chapter: relations
% Section: operations

\documentclass[../../../include/open-logic-section]{subfiles}

\begin{document}

\olfileid{sfr}{rel}{ops}
\olsection{Wat je met relaties kunt doen}

Soms willen we relaties veranderen of met elkaar combineren. In
\olref[sfr][rel][ord]{prop:stricttopartial} hebben we relaties \emph{verenigd}.
Daarvoor bekijken we twee relaties als verzamelingen van paren en voegen we die
verzamelingen samen. In \olref[sfr][rel][ord]{prop:partialtostrict} hebben we
op dezelfde manier het relatieve verschil van relaties bekeken. Nu volgen
enkele andere dingen die we met relaties kunnen doen.

\begin{defn}\ollabel{relationoperations} 
Neem twee relaties $R$ en $S$, en een willekeurige verzameling $A$.

De \emph{inverse relatie} van $R$ is $R^{-1} = \Setabs{\tuple{y, x}}{\tuple{x,
    y} \in R}$. Daarin is elk paar omgedraaid.

De \emph{samenstelling} van $R$ en $S$ is $(R \mid S) =
\{\tuple{x, z} : \exists y(Rxy \land Syz)\}$. Je gaat dus eerst via de
eerste relatie naar een element in het midden en daarna via de tweede relatie
naar het laatste element.

De \emph{restrictie} van $R$ tot $A$ is $\funrestrictionto{R}{A}= R
\cap A^2$. Je houdt daarmee alleen de paren uit die relatie over waarvan
beide delen in de opgegeven verzameling liggen.

Het door $R$ bepaalde \emph{beeld} van $A$ is $\funimage{R}{A} = \{y :
(\exists x \in A)Rxy\}$. Dit is de verzameling van alle tweede elementen die
via de relatie bij minstens één element uit de opgegeven verzameling horen.
\end{defn}

\begin{ex}
Neem de opvolgerrelatie $S \subseteq \Int^2$ op~$\Int$. Dat wil zeggen:
$S = \Setabs{\tuple{x, y} \in \Int^2}{x + 1 = y}$. Dus $Sxy$ geldt precies
als $x + 1 = y$.

$S^{-1}$ is de voorgangerrelatie op $\Int$, dus
$\Setabs{\tuple{x,y}\in\Int^2}{x -1 =y}$.

$S\mid S$ is
$ \Setabs{\tuple{x,y}\in\Int^2}{x + 2 =y}$

$\funrestrictionto{S}{\Nat}$ is de opvolgerrelatie op~$\Nat$.

$\funimage{S}{\{1,2,3\}}$ is $\{2, 3, 4\}$.
\end{ex}

\begin{defn}[Transitieve afsluiting]Neem een binaire relatie $R \subseteq A^2$.
	
De \emph{transitieve afsluiting} van~$R$ is $R^+ = \bigcup_{0 < n \in
\Nat} R^n$. De machten bouwen we stap voor stap op: $R^1 = R$ en
$R^{n+1} = R^n \mid R$.

De \emph{reflexief-transitieve afsluiting} van $R$ is $R^* = R^+ \cup
\Id{A}$. Daarmee voeg je ook alle paren toe van elementen uit de opgegeven
verzameling waarin het eerste en het tweede element hetzelfde zijn.
\end{defn}

\begin{ex}
Neem de opvolgerrelatie $S \subseteq \Int^2$. $S^2xy$ geldt precies als
$x + 2 = y$, $S^3xy$ precies als $x + 3 = y$, enzovoort. $S^+xy$ geldt dus
precies als $x + n = y$ voor een bepaalde $n \geq 1$. Of korter: $S^+xy$
geldt precies als $x < y$, en $S^*xy$ precies als $x \le y$.
\end{ex}

\begin{prob}
Laat zien dat de transitieve afsluiting van $R$ inderdaad transitief is.
\end{prob}
\end{document}
